\chapter{Generalization and Overfitting}

\section{Training and Test Performance}

In previous chapters, we formulated learning as minimizing empirical risk and studied how optimization algorithms find solutions. However, a central question remains:

\medskip

\noindent
\textbf{Why should a model that fits the training data well perform well on unseen data?}

\medskip

To address this, we distinguish between two types of performance.

\medskip

\noindent
\textbf{Training error.}
\[
\hat{R}_n(f) = \frac{1}{n} \sum_{i=1}^n \ell(f(x_i), y_i).
\]

\medskip

\noindent
\textbf{Test error (generalization error).}
\[
R(f) = \mathbb{E}_{(x,y)\sim P} \big[ \ell(f(x), y) \big].
\]

\medskip

\noindent
\textbf{Key distinction.}
\begin{itemize}
    \item Training error measures fit to observed data
    \item Test error measures performance on new data
\end{itemize}

\medskip

\noindent
A model generalizes well if its training error is close to its test error.

\medskip

\noindent
\textbf{Generalization gap.}
\[
R(f) - \hat{R}_n(f).
\]

\medskip

\noindent
\textbf{Insight.}  
Minimizing training error alone is not sufficient. The goal of learning is to minimize \emph{test error}, which is not directly observable.

\medskip

\noindent
\textbf{Practical approach.}  
We approximate test performance using a \emph{validation set}, separate from the training data.

\section{Overfitting and Underfitting}

The relationship between training and test error leads to two fundamental failure modes.

\medskip

\noindent
\textbf{Underfitting.}
\begin{itemize}
    \item Model is too simple to capture the structure of the data
    \item High training error
    \item High test error
\end{itemize}

\medskip

\noindent
\textbf{Overfitting.}
\begin{itemize}
    \item Model fits training data too closely, including noise
    \item Very low training error
    \item High test error
\end{itemize}

\medskip

\noindent
\textbf{Illustration.}  
A highly flexible model may interpolate all training points but behave erratically between them.

\medskip

\noindent
\textbf{Tradeoff.}  
As model complexity increases:
\begin{itemize}
    \item training error typically decreases,
    \item test error may decrease initially, then increase.
\end{itemize}

\medskip

\noindent
\textbf{Classical picture.}  
There exists an optimal level of complexity that balances underfitting and overfitting.

\medskip

\noindent
\textbf{Remark (modern perspective).}  
In deep learning, models often operate far beyond this classical regime, achieving low training error without necessarily suffering from poor generalization. This phenomenon will be revisited later.

\section{Model Capacity and Complexity}

To understand overfitting, we must quantify the expressive power of a model.

\medskip

\noindent
\textbf{Model capacity.}  
The ability of a hypothesis space $\mathcal{H}$ to fit a wide range of functions.

\medskip

\noindent
\textbf{Examples.}
\begin{itemize}
    \item Linear models: low capacity
    \item High-degree polynomials: high capacity
    \item Deep neural networks: very high capacity
\end{itemize}

\medskip

\noindent
\textbf{Qualitative effects of capacity.}
\begin{itemize}
    \item Low capacity $\rightarrow$ high bias, low variance
    \item High capacity $\rightarrow$ low bias, high variance
\end{itemize}

\medskip

\noindent
\textbf{Complexity measures (informal).}
\begin{itemize}
    \item Number of parameters
    \item Depth and width of networks
    \item Norms of parameters
\end{itemize}

\medskip

\noindent
\textbf{Important insight.}  
Capacity is not determined solely by the number of parameters. The structure of the model and the learning algorithm also play a crucial role.

\medskip

\noindent
\textbf{Example.}  
Two neural networks with the same number of parameters may behave very differently depending on:
\begin{itemize}
    \item architecture,
    \item initialization,
    \item optimization dynamics.
\end{itemize}

\medskip

\noindent
Thus, capacity should be viewed as a combination of representation, parameterization, and training procedure.

\section{Informal Generalization Principles}

Although a complete theory of generalization in modern machine learning is still evolving, several guiding principles are widely observed in practice.

\medskip

\noindent
\textbf{1. Simpler models tend to generalize better.}

\medskip

\noindent
This principle is often referred to as \emph{Occam's razor}: among models that fit the data, simpler ones are preferred.

\medskip

\noindent
Simplicity can take different forms:
\begin{itemize}
    \item fewer parameters,
    \item smaller norms,
    \item smoother functions.
\end{itemize}

\medskip

\noindent
\textbf{2. Regularization improves generalization.}

\medskip

\noindent
Regularization introduces additional constraints or penalties that favor simpler solutions.

\medskip

\noindent
Examples include:
\begin{itemize}
    \item weight decay,
    \item dropout,
    \item early stopping.
\end{itemize}

\medskip

\noindent
\textbf{3. More data improves generalization.}

\medskip

\noindent
As the sample size increases:
\begin{itemize}
    \item empirical risk better approximates true risk,
    \item estimation error decreases.
\end{itemize}

\medskip

\noindent
\textbf{4. Inductive bias is essential.}

\medskip

\noindent
Generalization depends on assumptions built into the model:
\begin{itemize}
    \item smoothness,
    \item invariance,
    \item compositional structure.
\end{itemize}

\medskip

\noindent
\textbf{5. Optimization influences generalization.}

\medskip

\noindent
The choice of optimization algorithm can affect which solution is found, even when multiple solutions achieve low training error.

\medskip

\noindent
This phenomenon is known as \emph{implicit regularization}.

\medskip

\noindent
\textbf{6. Noise can be beneficial.}

\medskip

\noindent
Stochasticity in training (e.g., SGD) can help avoid overfitting by preventing convergence to overly complex solutions.

\medskip

\noindent
\textbf{Summary insight.}  
Generalization is not determined by a single factor, but by the interaction between:
\begin{itemize}
    \item model capacity,
    \item data,
    \item inductive bias,
    \item optimization.
\end{itemize}

\section{A Unifying Perspective}

We can now reinterpret learning as the problem of balancing two objectives:
\begin{itemize}
    \item fitting the training data,
    \item maintaining the ability to generalize.
\end{itemize}

\medskip

\noindent
\textbf{Core tension.}  
\[
\text{Fit the data} \quad \text{vs} \quad \text{avoid overfitting}.
\]

\medskip

\noindent
\textbf{Practical strategy.}
\begin{enumerate}
    \item Choose an appropriate model class
    \item Control capacity through regularization
    \item Use validation to guide model selection
    \item Monitor training and test performance
\end{enumerate}

\medskip

\noindent
\textbf{Modern viewpoint.}  
Deep learning systems operate in regimes where classical intuitions are only partially valid, yet the underlying principles of inductive bias and structure remain essential.

\section{Outlook}

This chapter introduced generalization as a central challenge in learning:
\begin{itemize}
    \item distinguishing training and test performance,
    \item understanding overfitting and underfitting,
    \item relating model capacity to generalization.
\end{itemize}

\medskip

In the next chapter, we study regularization methods in greater detail, focusing on how to control model complexity in practice.

\medskip

\noindent
\textbf{Guiding principle.}  
A good model is not one that fits the data perfectly, but one that captures the underlying structure in a way that extends beyond the observed examples.